\section{Formalisation}
\label{sec:lean}

The repository contains a Lean~4 project \citep{demoura2021lean4} on Mathlib
\citep{mathlib2020} that states the conjecture and proves the results that
refute claim (L). A randomly initialised model is a parameter type with a
probability measure and an evaluation map, and one definition says that the
expected cost is strictly decreasing in the size of the model; the conjecture
and the learning claim are that definition for the nearest-point distance and
for a loss of the residual. Proved: \cref{thm:sandwich}, for inputs in any
seminormed group; \cref{thm:tube}, for an integrable distance in any
pseudo-metric space; \cref{thm:marginal}; and \cref{cor:phase}, from the
translation invariance of Haar measure on the circle. The last two combine into
the statement that the learning claim fails for sines of any amplitude, any
sequence of frequencies, any measurable loss and any measure on data points.
The project contains no \lean{sorry}, and the repository reports that every
theorem depends only on Lean's three standard axioms. Not formalised: claim (G)
itself, which is measured and not proved; \cref{lem:tube-area},
\cref{prop:length,prop:sine,prop:floor}; and \cref{cor:he}, so width enters the
formal development nowhere. \Cref{app:lean} lists the Lean names and versions.
